% !TEX root = ../main.tex
\chapter{术语表与符号总表}
\label{ch:glossary}

\section*{学习目标}
\begin{itemize}
  \item 把全书中最关键的术语、符号和层级区分压缩成可反复查阅的一份统一对照表。
  \item 避免在阅读 \cref{ch:bic,ch:dynamics,ch:workflow,ch:worked-example,ch:failure-modes} 等跨模块章节时发生语义漂移。
  \item 建立一套适合研究记录、仿真汇报和论文写作的统一命名习惯。
\end{itemize}

\section*{先修知识}
建议已通读全书，尤其是 \cref{ch:roadmap,ch:bic,ch:dynamics,ch:workflow,ch:worked-example,ch:failure-modes,ch:math-appendix,ch:device-appendix,ch:sim-appendix}。

\section*{本章路线图}
\ChapterModelLevel{跨层索引（L1→L4）}
本章先给出核心术语表，再给出符号总表，随后整理最常见的混淆对照，最后说明如何把本章真正用回工作流、例题与失败模式诊断。
为了帮助读者快速建立对全书符号体系的整体认知，\cref{fig:ch28_symbol_glossary}以树状图的形式展示了所有关键符号的分类架构。该图将数百个分散在各章的物理量系统地归纳为三大族群：基础物理（蓝色）、材料几何（绿色）和数值计算（橙色），并在每个类别下进一步细分为若干子类，标注出最具代表性的符号实例。建议初次接触本书时先将此图浏览一遍，后续遇到陌生符号时可迅速定位其所属范畴以减少翻阅时间。

\begin{figure}[htbp]
  \centering
  \begin{tikzpicture}[
      scale=0.82,
      transform shape,
      category/.style={rectangle, draw, rounded corners, fill=#1!20, minimum width=3.6cm, minimum height=0.85cm, align=center, font=\small},
      subcat/.style={rectangle, draw, dashed, fill=#1!10, minimum width=2.7cm, minimum height=0.65cm, align=center, font=\footnotesize},
      detail/.style={rectangle, draw, align=left, fill=white, text width=3.95cm, minimum height=1.7cm, font=\scriptsize, inner sep=4pt},
      wide/.style={rectangle, draw, rounded corners, align=left, fill=#1!8, text width=5.7cm, font=\scriptsize, inner sep=4pt},
      arrow/.style={->, >=stealth, thick}
    ]
    \node[category=blue] (physics) at (-5.1,5.0) {物理场量符号\\场、波矢、频率};
    \node[category=green] (materials) at (0,5.0) {材料与几何参数\\折射率、尺寸、填充};
    \node[category=orange] (numerics) at (5.1,5.0) {数值计算量\\离散步长、算符、本征值};

    \node[subcat=blue] (fields) at (-6.0,3.75) {电磁场\\$\mathbf{E},\mathbf{H},\mathbf{D},\mathbf{B}$};
    \node[subcat=blue] (waves) at (-4.2,3.75) {波动参数\\$\mathbf{k},\omega,\lambda,\beta$};
    \node[subcat=green] (optic) at (-0.9,3.75) {光学常数\\$n,n_g,\varepsilon,\chi^{(3)}$};
    \node[subcat=green] (geo) at (0.9,3.75) {几何尺寸\\$a,h,r,f,w,\Lambda$};
    \node[subcat=orange] (disc) at (4.2,3.75) {离散参数\\$N_x,N_t,\Delta x,\Delta t$};
    \node[subcat=orange] (ops) at (6.0,3.75) {算符与谱量\\$\hat{\Theta},\nabla\times,\lambda_n$};

    \foreach \src/\dst in {physics/fields,physics/waves,materials/optic,materials/geo,numerics/disc,numerics/ops}
      \draw[arrow] (\src) -- (\dst);

    \node[detail] at (-5.1,2.05) {
      \textbf{场量约定}\\
      $\mathbf{E}(\mathbf{r},t)=\mathrm{Re}[\tilde{\mathbf{E}}e^{-i\omega t}]$\\
      $\tilde{\mathbf{E}}$: 复振幅（相量）\\
      SI 制：$\mathbf{D}=\varepsilon_0\varepsilon_r\mathbf{E}$
    };

    \node[detail] at (0,2.05) {
      \textbf{波矢约定}\\
      $\mathbf{k}=k_x\hat{x}+k_y\hat{y}+k_z\hat{z}$\\
      $|\mathbf{k}|=n\omega/c=k_0n$\\
      布洛赫波矢通常约化到第一布里渊区
    };

    \node[detail] at (5.1,2.05) {
      \textbf{折射率类型}\\
      $n$: 相折射率\\
      $n_g=n-\lambda(\mathrm{d}n/\mathrm{d}\lambda)$: 群折射率\\
      $n_{\mathrm{eff}}=\beta/k_0$: 模式有效折射率
    };

    \node[wide=blue] at (-3.15,-0.25) {
      \textbf{常用希腊字母}\\
      $\alpha$: 损耗系数；$\gamma$: 约束因子；$\Gamma$: 光学限制因子\\
      $\kappa$: 耦合系数；$\eta$: 效率；$\lambda$: 真空波长；$\omega$: 角频率
    };

    \node[wide=green] at (3.15,-0.25) {
      \textbf{常见术语对照}\\
      PCSEL: 光子晶体面发射激光器；PC: 光子晶体；DBR: 分布布拉格反射镜\\
      CMT: 耦合模理论；PWE/FDTD/FEM: 常用数值求解方法
    };

    \node[rectangle, draw, dashed, rounded corners, align=center, fill=gray!5, text width=6.9cm, font=\scriptsize, inner sep=4pt] at (0,-2.25) {
      \textbf{量纲一致性检查}\\
      所有物理方程都应满足左右两侧量纲相同、求和项量纲一致。\\
      例：$\nabla\times\mathbf{E}$ 的单位为 $\mathrm{V/m^2}$。
    };
  \end{tikzpicture}
  \caption{本书主要符号与术语的分类索引。图中按场量、材料几何和数值变量三类汇总常见记号，并给出场量约定、波矢约定与量纲一致性提示。}
  \label{fig:ch28_symbol_glossary}
\end{figure}

\begin{IntuitionBox}
术语表不是为了把词翻译得“好看”，而是为了让证据链不失真。PCSEL 研究跨越开放电磁系统、半导体器件和多物理场建模，同一个词若在不同章节里含义漂移，结论就会在不知不觉中越级。
\end{IntuitionBox}

\section{如何使用本章}
本章最适合在三种情境下反复打开。第一，读到一个看似熟悉、但层级可能不同的词时，例如“threshold”“mode”“stability”。第二，准备把某个中间结果写进报告或论文时，用它检查自己是否把代理模型结论写成了器件结论。第三，团队讨论中出现“大家都在用同一个词，但好像在说不同东西”时，用本章统一语义。

如果需要一个最直接的用法，可以把本章与三张表配套使用：\cref{tab:workflow-example-failure,tab:worked-map,tab:worked-summary}。前两张表告诉你应该走哪条证据链，本章则保证链条上的每一个词和符号都指向同一层物理含义。

\section{核心术语表}
\begin{longtable}{@{}p{0.28\textwidth}p{0.66\textwidth}@{}}
\toprule
术语 & 教材化定义 \\
\midrule
\endhead
布洛赫模（Bloch mode） & 周期结构中满足布洛赫相位条件的本征场。它是无限周期分析的基本对象，但不自动等于有限器件中的工作模。 \\
$\Gamma$ 点模 & 布里渊区中心附近的模式家族。它与垂直辐射、偏振选择和面发射最直接相关，因此是 PCSEL 的核心候选模区域。 \\
带边模（band-edge mode） & 色散曲线局部斜率很小、群速度趋近于零附近的模式。它是被动色散概念，不自动等于最终激射模。 \\
被动共振（passive resonance） & 在无增益、无注入、无热回写条件下，由几何和材料分布决定的开放系统共振。 \\
有源阈值（active threshold） & 在给定增益、损耗、注入和工作点条件下，净增益刚好补偿总损耗的起振条件。 \\
阈值增益（threshold gain） & 线性光学层级下，为补偿模式总损耗所需的最小模增益。它是光学量，不等于器件阈值电流。 \\
阈值电流（threshold current） & 器件层级下，经过电注入、载流子输运、热回写和损耗共同决定的起振电流。 \\
导模共振（guided resonance） & 本来受波导约束、但因周期扰动耦合到辐射连续谱而可向外泄露的开放模式。 \\
连续谱束缚态（BIC） & 频率落在辐射连续谱中，却因对称性或干涉机制而不向外泄露的理想束缚态。 \\
对称性保护 BIC（symmetry-protected BIC） & 由于模式对称性与可辐射通道不匹配，导致耦合严格为零的 BIC。它常出现在 $\Gamma$ 点附近。 \\
偶然 BIC（accidental BIC） & 不是由严格对称性，而是由不同辐射通道之间的振幅和相位抵消形成的 BIC。 \\
准 BIC（quasi-BIC） & 当理想 BIC 因有限尺寸、扰动、无序或倾斜激发而获得有限泄露后形成的高 $Q$ 共振。它在真实 PCSEL 中更常见。 \\
辐射通道（radiation channel） & 模式可以向外界泄露能量的远场通道，包括不同偏振、角度和上下半空间的辐射分量。 \\
偏振选择 & 结构对不同偏振模式的辐射、损耗或增益施加不对称约束，从而导致某一偏振更容易成为稳定工作模。 \\
远场工程 & 通过改变晶格、扰动、有限阵列边界或纵向层栈，主动塑造主瓣、旁瓣和偏振纯度的设计行为。 \\
有效折射率近似 & 用等效二维折射率替代真实三维层状结构的代理模型。它适合快速筛选，但不能替代最终三维矢量验证。 \\
光学约束因子（optical confinement factor） & 模场与有源区或损耗区的重叠权重。它把材料层的性质映射到模层级。 \\
材料增益（material gain） & 材料本身在给定载流子密度和频率下提供的受激辐射增益。 \\
模增益（modal gain） & 材料增益经过模场重叠加权后作用于某一模式的有效增益。 \\
增益压缩（gain compression） & 随光子密度增大而导致微分增益下降的非线性效应。它直接进入动态响应与调制带宽分析。 \\
微分增益（differential gain） & 材料或模式增益对载流子密度的导数，是决定调制响应和弛豫振荡频率的重要量。 \\
Henry 因子（linewidth enhancement factor） & 载流子引起的折射率变化与增益变化之间的耦合因子，常记为 $\alpha_H$。它把振幅噪声耦合到相位噪声。 \\
弛豫振荡（relaxation oscillation） & 光子数与载流子数在平衡点附近的耦合振荡，是小信号调制响应中的核心特征频率。 \\
小信号调制响应 & 对工作点附近微小驱动扰动的线性频率响应，用于描述调制速度和动态稳定性。 \\
$3\,\mathrm{dB}$ 调制带宽 & 小信号调制响应幅度下降到低频值的 $1/\sqrt{2}$ 时对应的频率，常作为速度指标。 \\
鲁棒性窗口 & 在规定工艺误差、热漂移和注入偏差范围内，器件仍保持目标性能的参数窗口。 \\
容差分析 & 系统评估几何、材料、注入和热边界误差如何传播到模式排序、远场和器件性能的分析流程。 \\
验证闭环 & 从测量到参数提取、分层校准、模型预测再到回验的闭环过程，用于检验模型是否真正具备外推能力。 \\
回验（back-validation） & 使用未参与校准的观测量去检验模型预测能力的步骤。若只会拟合不会回验，模型就还不可靠。 \\
\bottomrule
\end{longtable}

\section{符号总表}
\begin{longtable}{@{}p{0.20\textwidth}p{0.24\textwidth}p{0.48\textwidth}@{}}
\toprule
符号 & 所属层级 & 含义 \\
\midrule
\endhead
$\omega,\tilde{\omega}$ & 电磁 & 实频角频率与复频率；复频率虚部表征损耗或增长。 \\
$Q$ & 电磁 & 品质因子，常由复频率或功率损耗定义。高 $Q$ 不自动等于最低阈值电流。 \\
$\alpha_\mathrm{rad}$ & 电磁 & 辐射损耗系数，反映模式向自由空间泄露的强弱。 \\
$\alpha_\mathrm{int}$ & 电磁/器件 & 内部损耗系数，包括自由载流子吸收、掺杂损耗、金属邻近损耗等。 \\
$\Gamma_\mathrm{opt}$ & 光学/器件 & 模场在有源区中的光学约束因子。 \\
$\Gamma_\mathrm{QW}$ & 器件 & 模场与量子阱区域的重叠因子。 \\
$\Gamma_\mathrm{loss}$ & 器件 & 模场与主要损耗区域的重叠因子。 \\
$g_\mathrm{mat}$ & 器件 & 材料增益。 \\
$g_\mathrm{mod}$ & 光学/器件 & 模增益，常近似写为 $\Gamma_\mathrm{opt} g_\mathrm{mat}$ 或更精细的空间积分形式。 \\
$g_\mathrm{th}$ & 光学 & 阈值增益，即补偿模式总损耗所需的最小模增益。 \\
$N(\rvec)$ & 器件 & 载流子密度空间分布。 \\
$J(\rvec)$ & 器件 & 电流密度空间分布。 \\
$T(\rvec)$ & 热/器件 & 温度空间分布。 \\
$\Delta n,\Delta g$ & 耦合模型 & 热或载流子回写到折射率和增益的变化量。 \\
$\tau_p$ & 动态学 & 光子寿命。 \\
$\tau_n$ & 动态学 & 有效载流子寿命。 \\
$S$ 或 $P$ & 动态学 & 腔内光子数或与之成比例的模式强度变量。 \\
$\beta_\mathrm{sp}$ & 动态学 & 自发辐射耦合因子。 \\
$\alpha_H$ & 动态学/噪声 & Henry 线宽增强因子。 \\
$\omega_R$ 或 $f_R$ & 动态学 & 弛豫振荡角频率或频率。 \\
$f_{3\mathrm{dB}}$ & 动态学 & 小信号调制响应的 $3\,\mathrm{dB}$ 带宽。 \\
$\varepsilon_g$ & 动态学 & 增益压缩系数。 \\
$\mu_n,\mu_p$ & 电学 & 电子与空穴迁移率。 \\
$D_n,D_p$ & 电学 & 电子与空穴扩散系数。 \\
$R_\mathrm{SRH}$ & 电学/器件 & Shockley--Read--Hall 复合率。 \\
$R_\mathrm{Auger}$ & 电学/器件 & Auger 复合率。 \\
$Q_\mathrm{Joule}$ & 热/器件 & Joule 热源项。 \\
$k_\mathrm{th}$ & 热学 & 热导率。 \\
$R_\mathrm{th}$ & 热学/实验接口 & 热阻，用于连接耗散功率与温升。 \\
$\kappa_1,\kappa_2$ & 半解析模型 & CWT 或玩具模型中的耦合系数。 \\
$\gamma_i,\gamma_r$ & 半解析模型 & 玩具模型中的内部损耗与辐射损耗。 \\
\bottomrule
\end{longtable}

\section{多义符号与上下文辨识}
\begin{longtable}{@{}p{0.18\textwidth}p{0.25\textwidth}p{0.49\textwidth}@{}}
\toprule
符号 & 常见语境 & 使用提醒 \\ \midrule
\endhead
$Q$ & 冷腔品质因子、辐射 $Q$、加载 $Q$ & 若不说明损耗来源，单写 $Q$ 很容易把辐射损耗、内部损耗和外部耦合混成同一量。 \\
$Q_c$ & 导带偏移系数 & 这是材料能带分配系数，不是品质因子 $Q$ 的下标变体；出现在外延和异质结章节时应优先按材料参数理解。 \\
$\Gamma$ & 布里渊区中心点、光学约束因子、量子阱重叠因子、某些降阶模型中的矩阵记号 & 章节切换时必须靠上下文或下标区分；$\Gamma$ 点与 $\Gamma_\mathrm{opt}$ 绝不是同一类对象，写成矩阵时也不应和约束因子混读。 \\
$\beta$ & slab 传播常数、自发辐射耦合因子 $\beta_\mathrm{sp}$ & 在器件和动态学章节尤其容易混淆，应尽量带下标或直接写全名。 \\
$g$ & 材料增益、模增益、阈值增益、微分增益 & 若没有下标，必须立即说明它属于材料层、模层还是导数层。 \\
$\alpha$ & 损耗系数、Henry 因子 $\alpha_H$ & 前者是损耗量，后者是线宽增强因子；两者物理意义完全不同。 \\
$N$ & 载流子密度、样本数、阵列周期数 & 数值章节和器件章节中应避免只写孤立的 $N$ 而不加说明。 \\
\bottomrule
\end{longtable}

\begin{ExampleBox}
\textbf{使用建议。} 遇到多义符号时，先看它第一次出现在哪一章；再看它是否带下标；最后判断它属于电磁、器件、动态学还是工作流语境。对全书最常见的几个多义符号，最稳妥的写法是不要让 $Q$、$\Gamma$、$\beta$、$g$ 脱离上下文孤立出现。
\end{ExampleBox}

\begin{RigorousBox}
全书中凡出现“threshold”“stability”“mode selection”这类词，默认都必须带层级。若没有说明它属于被动电磁层级、简化器件代理层级，还是器件级工作点层级，那么表述就是不完整的。
\end{RigorousBox}

\section{高频混淆对照}
\begin{longtable}{@{}p{0.30\textwidth}p{0.62\textwidth}@{}}
\toprule
易混淆对 & 教材化区分句 \\
\midrule
\endhead
带边模 vs 激射模 & 前者是被动色散对象，后者是增益、损耗、注入和热回写共同决定的工作结果。 \\
被动共振 vs 有源阈值 & 前者描述结构本身的开放电磁行为，后者描述增益补偿后的起振条件。 \\
阈值增益 vs 阈值电流 & 前者是光学量，后者是器件量；前者可以来自冷腔和增益代理，后者必须经过电热回写。 \\
BIC vs 高 $Q$ 共振 & BIC 是理想耦合严格为零的束缚态，高 $Q$ 共振只说明泄露很小，不说明耦合严格为零。 \\
准 BIC vs 低阈值模 & 准 BIC 可能提供较低辐射损耗，但若注入、热和偏振选择不利，仍不一定成为最低阈值器件模。 \\
亮模 vs 先激射模 & 亮模说明向外辐射强，不说明净阈值最低。亮模甚至可能因为辐射损耗大而更晚起振。 \\
无限周期 vs 有限器件 & 前者给候选模式池和局部色散信息，后者决定边缘泄露、远场、尺寸收敛和真实工作模。 \\
均匀增益 vs 电注入增益 & 前者是代理模型，后者由 $J(\rvec)$、$N(\rvec)$ 和温度场共同决定。 \\
冷腔频率 vs 工作频率 & 工作频率还会受载流子、温升和非线性效应拖动。 \\
校准量 vs 回验量 & 校准量用于拟合模型参数，回验量用于检验模型外推能力；两者不能混用。 \\
名义值最优 vs 鲁棒性最优 & 前者只在单点参数上最好，后者要求在误差和漂移存在时仍保持性能窗口。 \\
仿真一致 vs 物理正确 & 两个求解器给出相近数字只是必要条件，不是充分条件；仍需检查边界、层级和守恒关系。 \\
\bottomrule
\end{longtable}

\begin{ExampleBox}
合规写法示例：“该结构在被动开放电磁模型下显示出接近准 BIC 的高 $Q$ 候选模，因此在光学层级上具有较低辐射损耗潜力；其器件级阈值电流与线宽仍需结合 \cref{ch:dynamics,ch:device-appendix} 中的动态学、注入与热回写模型评估。”
\end{ExampleBox}

\begin{PitfallBox}
把“存在一个高 $Q$ 的 $\Gamma$ 点候选模”直接改写成“器件一定会以单模低阈值工作”，是全书最需要避免的术语越级。类似地，把“调制响应更快”写成“线宽一定更窄”也是把动态学结论和噪声结论混在了一起。
\end{PitfallBox}

\section{如何把本章用回工作流、例题与失败模式}
\cref{ch:workflow} 负责说明建模链应该怎样向前推进；\cref{ch:worked-example} 把这条链真正走了一遍；\cref{ch:failure-modes} 则把同一条链倒过来，检查哪里最容易断。本章的作用，是确保这三条阅读路径使用的是同一套词汇和符号。

更具体地说：
\begin{itemize}
  \item 当你在 \cref{ch:bic} 里看到 BIC、quasi-BIC、radiation channel 时，应回到本章确认它们与高 $Q$、亮模、远场工程的边界。
  \item 当你在 \cref{ch:dynamics} 里看到 $\alpha_H$、$\omega_R$、$f_{3\mathrm{dB}}$ 时，应回到本章确认这些量属于动态学和噪声层级，而不是静态阈值层级。
  \item 当你在 \cref{ch:worked-example} 里整理中间结果时，可直接用本章检查每一个输出量属于“候选模证据”“有限尺寸证据”还是“器件代理证据”。
  \item 当你在 \cref{ch:failure-modes} 中排错时，可用本章快速判断自己究竟是符号混用、术语越级，还是模型层级本身已经错位。
\end{itemize}

\section*{本章小结}
本章把全书中最容易漂移的术语和符号压缩成一份统一工具表。真正有效的用法，不是偶尔查一个词，而是在写公式、做工作流、整理例题中间结果和诊断失败模式时，持续用它检查自己是否发生了层级混用、概念越级或符号漂移。

\section*{练习题}
\begin{enumerate}
  \item \basicq 从 \cref{ch:bic,ch:dynamics} 中各找出两个你最容易混淆的术语，分别写出一句不会越级的定义。

  \item \basicq 用本章的混淆对照表，重写一句你在研究笔记里曾经写过的“高 $Q$ 因此低阈值”式表述。

  \item \midq 以 \cref{ch:worked-example} 的一张中间结果表为对象，给出每个中间量的层级标签：被动光学、有限尺寸、器件代理或动态学。

  \item \hardq 说明为什么在实验校准中，回验量不能再反过来参与同一轮参数拟合。
\end{enumerate}

\section*{建议继续阅读}
\begin{itemize}
  \item 若你在整理理论与辐射主线，建议与 \cref{ch:bic,ch:math-appendix} 对照阅读。 这条阅读的重点是把本章的输出顺着主线接到后续模块，而不是停成孤立知识点。

  \item 若你在整理调制、线宽与噪声主线，建议与 \cref{ch:dynamics,ch:device-appendix} 对照阅读。 这条阅读的重点是把本章的输出顺着主线接到后续模块，而不是停成孤立知识点。

  \item 若你在整理项目执行链，建议与 \cref{ch:workflow,ch:worked-example,ch:failure-modes,ch:sim-appendix} 对照阅读。 这条阅读的重点是把本章的输出顺着主线接到后续模块，而不是停成孤立知识点。
\end{itemize}


